```latex
\chapter{KẾT LUẬN VÀ KIẾN NGHỊ}

\section{Kết luận}

Trong bối cảnh chuyển đổi số đang diễn ra mạnh mẽ trong ngành ngân hàng, việc ứng dụng trí tuệ nhân tạo vào hoạt động quản trị rủi ro tín dụng ngày càng trở thành một xu hướng tất yếu. Xếp hạng tín dụng không chỉ đóng vai trò hỗ trợ quyết định cấp tín dụng mà còn là công cụ quan trọng giúp ngân hàng kiểm soát rủi ro, nâng cao chất lượng danh mục tín dụng và đảm bảo an toàn hoạt động.

Xuất phát từ yêu cầu thực tiễn đó, luận văn đã thực hiện nghiên cứu đề tài:

\begin{center}
\textit{"Ứng dụng trí tuệ nhân tạo trong xếp hạng tín dụng: nghiên cứu từ một ngân hàng thương mại ở Việt Nam"}
\end{center}

Mục tiêu của nghiên cứu là xây dựng và đánh giá các mô hình trí tuệ nhân tạo nhằm hỗ trợ bài toán xếp hạng tín dụng khách hàng dựa trên dữ liệu thực tế của ngân hàng thương mại.

Thông qua quá trình nghiên cứu, luận văn đã hoàn thành các nội dung chính bao gồm:

\begin{itemize}
    \item Tổng quan cơ sở lý thuyết về tín dụng, rủi ro tín dụng, xếp hạng tín dụng và trí tuệ nhân tạo trong lĩnh vực tài chính ngân hàng.
    
    \item Tổng hợp và phân tích các nghiên cứu trong nước và quốc tế liên quan đến bài toán chấm điểm tín dụng và xếp hạng tín dụng.
    
    \item Xây dựng quy trình xử lý dữ liệu tín dụng bao gồm làm sạch dữ liệu, xử lý dữ liệu thiếu, phát hiện dữ liệu bất thường, xây dựng đặc trưng và xử lý mất cân bằng dữ liệu.
    
    \item Huấn luyện và đánh giá nhiều mô hình học máy khác nhau như Logistic Regression, Decision Tree, Random Forest, XGBoost, LightGBM và CatBoost.
    
    \item So sánh hiệu quả của các mô hình trên bộ dữ liệu thực tế và lựa chọn mô hình có hiệu năng tốt nhất.
    
    \item Ứng dụng phương pháp SHAP nhằm giải thích kết quả dự báo của mô hình, nâng cao tính minh bạch trong quá trình ra quyết định tín dụng.
\end{itemize}

Kết quả nghiên cứu cho thấy các mô hình học máy hiện đại có khả năng hỗ trợ hiệu quả cho bài toán xếp hạng tín dụng. So với các phương pháp truyền thống, các mô hình Ensemble Learning như Random Forest, XGBoost, LightGBM và CatBoost đạt hiệu năng dự báo cao hơn, đồng thời có khả năng khai thác các quan hệ phi tuyến trong dữ liệu tín dụng.

Bên cạnh đó, việc sử dụng các kỹ thuật xử lý mất cân bằng dữ liệu như SMOTE giúp cải thiện khả năng nhận diện các nhóm khách hàng có mức độ rủi ro cao, vốn thường chiếm tỷ lệ nhỏ trong tập dữ liệu.

Kết quả nghiên cứu cũng cho thấy các yếu tố như lịch sử tín dụng, số ngày quá hạn, dư nợ hiện tại, số lượng khoản vay đang hoạt động và thời gian quan hệ tín dụng là những yếu tố quan trọng ảnh hưởng đến kết quả xếp hạng tín dụng của khách hàng.

\section{Đóng góp của nghiên cứu}

\subsection{Đóng góp về mặt khoa học}

Về mặt học thuật, luận văn đã hệ thống hóa các cơ sở lý thuyết liên quan đến:

\begin{itemize}
    \item Rủi ro tín dụng;
    \item Xếp hạng tín dụng;
    \item Trí tuệ nhân tạo;
    \item Học máy trong lĩnh vực tài chính ngân hàng.
\end{itemize}

Nghiên cứu đồng thời cung cấp thêm bằng chứng thực nghiệm về hiệu quả của các thuật toán học máy hiện đại trong bài toán xếp hạng tín dụng tại Việt Nam.

Kết quả nghiên cứu góp phần bổ sung nguồn tài liệu tham khảo cho các nghiên cứu tiếp theo về:

\begin{itemize}
    \item Credit Scoring;
    \item Credit Rating;
    \item Default Prediction;
    \item Early Warning System;
    \item Explainable Artificial Intelligence.
\end{itemize}

\subsection{Đóng góp về mặt thực tiễn}

Về mặt thực tiễn, nghiên cứu đã xây dựng được quy trình hoàn chỉnh từ xử lý dữ liệu đến xây dựng mô hình xếp hạng tín dụng có thể áp dụng trong môi trường ngân hàng.

Mô hình đề xuất có thể hỗ trợ:

\begin{itemize}
    \item Đánh giá mức độ tín nhiệm của khách hàng.
    \item Cảnh báo sớm rủi ro tín dụng.
    \item Hỗ trợ cán bộ tín dụng trong quá trình ra quyết định.
    \item Nâng cao chất lượng danh mục tín dụng.
    \item Giảm thiểu rủi ro phát sinh nợ xấu.
\end{itemize}

Ngoài ra, kết quả nghiên cứu có thể được tích hợp vào hệ thống xếp hạng tín dụng nội bộ hoặc hệ thống hỗ trợ quyết định của ngân hàng.

\section{Hạn chế của nghiên cứu}

Mặc dù đạt được các kết quả tích cực, nghiên cứu vẫn tồn tại một số hạn chế nhất định.

Thứ nhất, dữ liệu nghiên cứu chỉ được thu thập từ một ngân hàng thương mại. Do đó, kết quả nghiên cứu có thể chưa phản ánh đầy đủ đặc điểm của toàn bộ hệ thống ngân hàng tại Việt Nam.

Thứ hai, bộ dữ liệu nghiên cứu chủ yếu bao gồm các thông tin tín dụng truyền thống. Các nguồn dữ liệu thay thế như dữ liệu giao dịch số, dữ liệu hành vi khách hàng, dữ liệu mạng xã hội hoặc dữ liệu viễn thông chưa được khai thác.

Thứ ba, nghiên cứu tập trung vào các thuật toán học máy truyền thống và các mô hình boosting. Các mô hình học sâu (Deep Learning) chưa được xem xét trong phạm vi nghiên cứu.

Thứ tư, mô hình được đánh giá chủ yếu trên dữ liệu lịch sử. Việc triển khai và đánh giá hiệu quả trong môi trường thực tế chưa được thực hiện.

Thứ năm, mặc dù đã sử dụng SHAP để giải thích mô hình, khả năng giải thích của các mô hình học máy phức tạp vẫn còn là một thách thức trong lĩnh vực tài chính ngân hàng.

\section{Kiến nghị đối với ngân hàng thương mại}

Từ kết quả nghiên cứu, tác giả đề xuất một số kiến nghị đối với ngân hàng thương mại như sau:

\subsection{Nâng cao chất lượng dữ liệu}

Ngân hàng cần tiếp tục đầu tư vào hệ thống quản trị dữ liệu nhằm:

\begin{itemize}
    \item Chuẩn hóa dữ liệu tín dụng.
    \item Nâng cao chất lượng dữ liệu đầu vào.
    \item Giảm thiểu dữ liệu thiếu và dữ liệu sai lệch.
    \item Xây dựng kho dữ liệu tập trung.
\end{itemize}

\subsection{Ứng dụng trí tuệ nhân tạo trong xếp hạng tín dụng}

Ngân hàng nên từng bước tích hợp các mô hình trí tuệ nhân tạo vào hệ thống xếp hạng tín dụng hiện có nhằm:

\begin{itemize}
    \item Tăng độ chính xác trong đánh giá tín dụng.
    \item Rút ngắn thời gian phê duyệt hồ sơ.
    \item Tăng khả năng cảnh báo sớm rủi ro.
    \item Hỗ trợ tự động hóa quy trình tín dụng.
\end{itemize}

\subsection{Phát triển hệ thống Explainable AI}

Trong môi trường ngân hàng, khả năng giải thích quyết định là yếu tố đặc biệt quan trọng.

Do đó, ngân hàng cần:

\begin{itemize}
    \item Xây dựng các mô hình có khả năng giải thích.
    \item Kết hợp SHAP hoặc các công cụ Explainable AI.
    \item Đảm bảo tính minh bạch trong quá trình ra quyết định.
\end{itemize}

\section{Hướng phát triển trong tương lai}

Trong tương lai, nghiên cứu có thể được mở rộng theo một số hướng sau:

\subsection{Mở rộng quy mô dữ liệu}

Thu thập dữ liệu từ nhiều ngân hàng khác nhau nhằm xây dựng mô hình có khả năng tổng quát hóa tốt hơn.

\subsection{Khai thác dữ liệu phi truyền thống}

Bổ sung các nguồn dữ liệu thay thế như:

\begin{itemize}
    \item Dữ liệu giao dịch số.
    \item Dữ liệu thanh toán điện tử.
    \item Dữ liệu hành vi khách hàng.
    \item Dữ liệu từ ví điện tử.
\end{itemize}

\subsection{Ứng dụng Deep Learning}

Nghiên cứu các mô hình học sâu như:

\begin{itemize}
    \item Artificial Neural Network (ANN);
    \item Deep Neural Network (DNN);
    \item AutoEncoder;
    \item Transformer;
    \item TabNet.
\end{itemize}

để cải thiện hiệu quả xếp hạng tín dụng.

\subsection{Xây dựng hệ thống cảnh báo sớm}

Kết hợp mô hình xếp hạng tín dụng với hệ thống cảnh báo sớm nhằm dự báo khả năng chuyển nhóm nợ của khách hàng trong tương lai.

\subsection{Triển khai thực tế}

Nghiên cứu khả năng triển khai mô hình trên nền tảng Web hoặc Mobile Banking nhằm hỗ trợ cán bộ tín dụng trong môi trường thực tế.

\section{Kết luận chương}

Chương này đã tổng kết toàn bộ nội dung nghiên cứu của luận văn, làm rõ những kết quả đạt được, các đóng góp khoa học và thực tiễn của đề tài.

Kết quả nghiên cứu cho thấy trí tuệ nhân tạo có tiềm năng lớn trong việc nâng cao hiệu quả xếp hạng tín dụng tại các ngân hàng thương mại Việt Nam. Các mô hình học máy hiện đại không chỉ giúp cải thiện độ chính xác dự báo mà còn hỗ trợ ngân hàng trong công tác quản trị rủi ro tín dụng và chuyển đổi số.

Mặc dù vẫn còn một số hạn chế, nghiên cứu đã tạo nền tảng cho các hướng phát triển tiếp theo trong lĩnh vực ứng dụng trí tuệ nhân tạo vào quản trị rủi ro tín dụng tại Việt Nam.
```
